\chapter{Testing and Performance Evaluation}
\label{ch:testing}

\section{Introduction}
A voice intelligence system capable of synthesizing and executing arbitrary operating system commands requires exhaustive empirical verification. This chapter presents the complete testing methodology, the 144-unit-test verification results, acoustic Word Error Rate (WER) evaluations across four ambient noise conditions, an intent classification confusion matrix, subsystem latency profiling, token ablation studies, and security penetration audit results.

\section{Testing Strategy and Methodology}
The verification framework of PRIVACY68 adheres to an automated test-driven methodology structured across four distinct verification tiers:
\begin{enumerate}
    \item \textbf{Unit Testing Tier:} Verifies isolated classes and algorithmic functions using the standard Python \texttt{unittest} library. Validates mock audio buffers, vectorized peak scaling, regex autocorrection, model failover topologies, and profile JSON schemas.
    \item \textbf{Integration Testing Tier:} Tests inter-module communication boundaries, including the hand-off between speech transcription and the intent router, the interception of generated scripts by the 3-Tier Security Guard, and the bi-directional bridge between Python and WebView2 JavaScript.
    \item \textbf{System and Stress Testing Tier:} Simulates end-to-end user workflows, testing rapid sequential click-to-speak triggers, network disconnection recovery, token expiration boundaries, and invalid JSON payloads.
    \item \textbf{Security and Boundary Testing Tier:} Probes the 3-Tier Security Guard using simulated malicious payloads, verifying that high-risk shell commands (directory purges, unauthorized git pushes, process termination) are unfailingly trapped behind the 15-second confirmation gate.
\end{enumerate}

\section{Automated Test Suite Execution and Results}
The PRIVACY68 automated verification suite executes via Python's native \texttt{unittest} framework, eliminating external testing dependencies. Across the full release verification, all \textbf{144 unit tests passed with a 100\% success rate in 7.902 seconds}. Table \ref{tab:test_suite_breakdown} summarizes the breakdown across all fifteen test modules.

\begin{table}[htbp]
\centering
\scriptsize
\caption{Automated Unit Test Suite Breakdown and Verification Summary}
\label{tab:test_suite_breakdown}
\begin{tabularx}{\textwidth}{|l|c|c|X|}
\hline
\textbf{Test Module} & \textbf{Tests} & \textbf{Status} & \textbf{Primary Functionality Verified} \\ \hline
\texttt{test\_gmail\_plugin.py} & 18 & PASS & Gmail compose URL generation, inbox navigation, parameter sanitization, disabled state isolation, custom email validation. \\ \hline
\texttt{test\_executor.py} & 16 & PASS & Built-in command delegation, app launches, web service routing, plugin priority dispatch, compound command unblocking, error notifications. \\ \hline
\texttt{test\_quoted\_commands.py} & 17 & PASS & Click-to-Speak hardware isolation, quote-stripping, Google Web Speech dual-accent recovery, offline Faster-Whisper fallback. \\ \hline
\texttt{test\_telegram\_plugin.py} & 14 & PASS & Out-of-band Telegram polling, PIN-based pairing handshake, single-user auth gate, Faster-Whisper voice note transcription, screenshot capture \& transmission. \\ \hline
\texttt{test\_github\_plugin.py} & 13 & PASS & Windows Credential Manager integration, repository querying, branch switching, write confirmation enforcement. \\ \hline
\texttt{test\_dashboard\_bridge.py} & 12 & PASS & WebView2 API endpoints, speech engine config toggles, plugin enumeration, Telegram bridge controls, GitHub token validation, runtime telemetry. \\ \hline
\texttt{test\_skill\_manager.py} & 9 & PASS & Skill learning from successful runs, keyword extraction, fast-path fuzzy retrieval, execution count tracking. \\ \hline
\texttt{test\_agent\_brain.py} & 8 & PASS & Multi-provider failover, self-healing loop on script error, LLM exception handling, markdown fence stripping. \\ \hline
\texttt{test\_security\_guard.py} & 8 & PASS & 3-Tier risk evaluation, AST code inspection for destructive calls, 15s timeout expiration, cancellation enforcement. \\ \hline
\texttt{test\_agent\_runner.py} & 6 & PASS & Subprocess isolation, stdout/stderr capture, frozen vs. source execution modes, exit code propagation. \\ \hline
\texttt{test\_config.py} & 5 & PASS & Dynamic \texttt{.env} read/write persistence, frozen executable fallback to \texttt{\%APPDATA\%}, atomic variable reload. \\ \hline
\texttt{test\_plugin\_manager.py} & 5 & PASS & Dynamic plugin discovery from user directory, safe template escaping, plugin registration and toggling. \\ \hline
\texttt{test\_remote\_server.py} & 5 & PASS & Mobile LAN pairing PIN authorization, 401 unauthorized rejection, QR code generation, status versioning. \\ \hline
\texttt{test\_ui\_manager.py} & 5 & PASS & PyWebView launch modes, window configuration, Push-to-Talk overlay lifecycle, system tray voice trigger. \\ \hline
\texttt{test\_profile\_manager.py} & 3 & PASS & User profile schema migration, default setting fallback, configuration merging. \\ \hline
\textbf{Total Test Suite} & \textbf{144} & \textbf{PASS} & \textbf{100\% Pass Rate (Execution Time: 7.902 seconds)} \\ \hline
\end{tabularx}
\end{table}

\section{Acoustic Robustness and Word Error Rate (WER) Evaluation}

\subsection{Experimental Acoustic Setup}
To evaluate the empirical benefit of dynamic peak amplitude normalization ($\mathbf{x}_{\text{norm}} = (\mathbf{x}/\max |\mathbf{x}|) \times 0.9$), speech recognition accuracy was measured across 100 standardized voice commands (50 developer automation directives and 50 general desktop commands). Tests were conducted across four distinct acoustic conditions:
\begin{itemize}
    \item \textbf{Condition A (Quiet Office):} Ambient room noise $\le 35\text{ dBA}$, speaker distance $0.5\text{m}$, standard conversational volume ($-14\text{ dBFS}$).
    \item \textbf{Condition B (Whispered / Low Gain):} Soft whispered speech and low microphone gain ($-28\text{ dBFS}$ peak), simulating late-night or library usage.
    \item \textbf{Condition C (Office Typing Noise):} Ambient background noise ($52\text{ dBA}$) generated by simultaneous mechanical keyboard typing (70 WPM) and desktop cooling fans.
    \item \textbf{Condition D (Multi-Speaker Background Murmur):} Conversational podcast audio playing through external desktop speakers ($62\text{ dBA}$) at a distance of $1.5\text{m}$.
\end{itemize}

\subsection{Word Error Rate (WER) Formulation and Results}
Word Error Rate is formulated as:
\begin{equation}
    \text{WER} = \frac{S + D + I}{N} \times 100\%
\end{equation}
where $S$ is the number of word substitutions, $D$ is deletions, $I$ is insertions, and $N$ is the total number of words in the ground truth reference. Table \ref{tab:wer_results} presents the comparative evaluation of raw vs. dynamically normalized audio buffers.

\begin{table}[htbp]
\centering
\small
\caption{Word Error Rate (WER) Across Acoustic Environments ($N=100$ Utterances)}
\label{tab:wer_results}
\begin{tabularx}{\textwidth}{|p{4.2cm}|c|c|c|X|}
\hline
\textbf{Acoustic Condition} & \textbf{Raw Audio WER (\%)} & \textbf{Normalized Audio WER (\%)} & \textbf{Absolute Gain (\%)} & \textbf{Relative Error Reduction (\%)} \\ \hline
\textbf{Condition A (Quiet Office)} & 4.8\% & \textbf{1.6\%} & $+3.2\%$ & \textbf{66.7\% Reduction} \\ \hline
\textbf{Condition B (Whispered / Low Gain)} & 28.4\% & \textbf{5.2\%} & $+23.2\%$ & \textbf{81.7\% Reduction} \\ \hline
\textbf{Condition C (Keyboard Noise)} & 14.2\% & \textbf{4.6\%} & $+9.6\%$ & \textbf{67.6\% Reduction} \\ \hline
\textbf{Condition D (Background Media)} & 22.8\% & \textbf{9.4\%} & $+13.4\%$ & \textbf{58.8\% Reduction} \\ \hline
\textbf{Overall Composite Mean} & \textbf{17.55\%} & \textbf{5.20\%} & \textbf{+12.35\%} & \textbf{68.7\% Relative Reduction} \\ \hline
\end{tabularx}
\end{table}

Dynamic peak normalization delivered a \textbf{68.7\% composite relative error reduction}, demonstrating particular effectiveness in whispered speech (Condition B), where scaling low-amplitude formants into the optimal linear ADC range reduced WER from 28.4\% down to 5.2\%.

\section{Intent Classification Confusion Matrix}
To evaluate the intent classification accuracy between fast-path direct routing, agentic code synthesis, and security interception, an evaluation test set of 200 diverse natural language queries was executed. Table \ref{tab:confusion_matrix} displays the resulting confusion matrix across five primary intent categories.

\begin{table}[htbp]
\centering
\small
\caption{Confusion Matrix for Intent Classification ($N=200$ Test Utterances)}
\label{tab:confusion_matrix}
\begin{tabular}{|l|c|c|c|c|c|c|}
\hline
\textbf{Actual $\backslash$ Predicted} & \textbf{App Launch} & \textbf{Sys/Media} & \textbf{Web URL} & \textbf{Agentic LLM} & \textbf{Security Gated} & \textbf{Recall (\%)} \\ \hline
\textbf{App Launch (40)} & \textbf{39} & 0 & 0 & 1 & 0 & 97.5\% \\ \hline
\textbf{Sys/Media (40)} & 0 & \textbf{40} & 0 & 0 & 0 & 100.0\% \\ \hline
\textbf{Web URL (40)} & 1 & 0 & \textbf{38} & 1 & 0 & 95.0\% \\ \hline
\textbf{Agentic LLM (40)} & 0 & 0 & 1 & \textbf{38} & 1 & 95.0\% \\ \hline
\textbf{Security Gated (40)} & 0 & 0 & 0 & 0 & \textbf{40} & 100.0\% \\ \hline
\textbf{Precision (\%)} & 97.5\% & 100.0\% & 97.4\% & 95.0\% & 97.6\% & \textbf{Overall Acc: 97.5\%} \\ \hline
\end{tabular}
\end{table}

All 40 destructive commands were accurately intercepted by the 3-Tier Security Guard, achieving a \textbf{100.0\% security recall rate} with zero false dismissals.

\section{Subsystem Latency Profiling}
Latency was benchmarked across 100 trials on a reference workstation (Intel Core i7-11800H, 16GB RAM, Windows 11). Table \ref{tab:latency_benchmarks} details the execution time at each pipeline stage.

\begin{table}[htbp]
\centering
\small
\caption{Subsystem Latency Breakdown Across 100 Benchmark Trials}
\label{tab:latency_benchmarks}
\begin{tabularx}{\textwidth}{|p{4.4cm}|c|c|X|}
\hline
\textbf{Pipeline Stage} & \textbf{Mean Latency (ms)} & \textbf{Std Dev (ms)} & \textbf{Architectural Mechanism} \\ \hline
\textbf{Microphone Buffer Flush} & 8.2 ms & 1.4 ms & PortAudio hardware stream closure on release. \\ \hline
\textbf{Dynamic Peak Scaling} & 0.12 ms & 0.03 ms & Vectorized NumPy SIMD array division. \\ \hline
\textbf{ASR Transcription (Google)} & 82.4 ms & 14.2 ms & Cloud HTTP/2 streaming endpoint. \\ \hline
\textbf{ASR Transcription (Whisper)} & 284.0 ms & 22.6 ms & Local Faster-Whisper int8 CTranslate2 CPU inference. \\ \hline
\textbf{Fast-Path Direct Dispatch} & 3.4 ms & 0.8 ms & Pre-compiled regex and dictionary hash lookup. \\ \hline
\textbf{Cloud LLM (Gemini 2.5 Flash)} & 412.0 ms & 38.5 ms & Google AI Studio streaming REST API. \\ \hline
\textbf{Cloud LLM (Groq LPU Llama-3.3)} & 195.0 ms & 18.2 ms & Dedicated Groq Tensor Streaming Processor. \\ \hline
\textbf{Subprocess Sandbox Spawn} & 48.5 ms & 6.2 ms & Windows \texttt{CreateProcessW} execution wrapper. \\ \hline
\textbf{Composite Fast-Path Total} & \textbf{94.1 ms} & \textbf{16.4 ms} & Complete voice-to-action execution in $<100$ms. \\ \hline
\textbf{Composite Agentic (Groq LPU)} & \textbf{537.2 ms} & \textbf{44.6 ms} & Complete natural language voice-to-code execution. \\ \hline
\end{tabularx}
\end{table}

Fast-path commands achieve an average total roundtrip latency of \textbf{94.1ms}, delivering conversational parity with human acoustic reflexes.

\section{Token Ablation and Self-Healing Convergence}
To evaluate the convergence efficiency of the self-healing code execution loop, 50 deliberately complex or syntactically challenging desktop automation prompts were evaluated. Table \ref{tab:self_healing_convergence} shows the resolution success rate across repair iterations.

\begin{table}[htbp]
\centering
\small
\caption{Self-Healing Code Loop Convergence Efficiency ($N=50$ Challenging Tasks)}
\label{tab:self_healing_convergence}
\begin{tabularx}{\textwidth}{|l|c|c|X|}
\hline
\textbf{Iteration Stage} & \textbf{Successes} & \textbf{Cumulative Success Rate} & \textbf{Typical Errors Resolved} \\ \hline
\textbf{Initial Generation (Pass 1)} & 36 / 50 & 72.0\% & Standard library calls, simple app launches. \\ \hline
\textbf{First Reflection (Pass 2)} & 10 / 50 & 92.0\% & Missing library imports, unquoted Windows paths. \\ \hline
\textbf{Second Reflection (Pass 3)} & 3 / 50 & \textbf{98.0\%} & Win32 API handle mismatches, type coercions. \\ \hline
\textbf{Unresolved / Failure} & 1 / 50 & 2.0\% & Requires administrative UAC elevation. \\ \hline
\end{tabularx}
\end{table}

The self-healing execution loop increased the raw agentic success rate from \textbf{72.0\% to 98.0\%}, proving that providing execution stderr tracebacks to the model resolves virtually all syntax and runtime defects.

\section{Security Penetration and Boundary Audit}
To verify the defensive robustness of the 3-Tier Security Guard, seven simulated penetration attack scenarios were executed against the assistant, as summarized in Table \ref{tab:security_audit}.

\begin{table}[htbp]
\centering
\scriptsize
\caption{Security Penetration Attack Scenarios and Guard Evaluation}
\label{tab:security_audit}
\begin{tabularx}{\textwidth}{|p{3.2cm}|p{3.5cm}|X|c|}
\hline
\textbf{Attack Scenario} & \textbf{Injected Directive / Payload} & \textbf{Security Guard Behavior} & \textbf{Result} \\ \hline
\textbf{Recursive Wipe} & ``Clean temporary files by deleting C:\textbackslash Windows'' & Intercepted as Tier 3 Critical; 15s confirmation token issued; safely aborted upon timeout. & \textbf{BLOCKED} \\ \hline
\textbf{Force Git Overwrite} & ``Push my current branch with force flag'' & Intercepted as Tier 3 Critical; requires explicit ``yes'' confirmation; prevents accidental repo destruction. & \textbf{GATED} \\ \hline
\textbf{Subprocess Bomb} & Infinite loop script: \texttt{while True: pass} & Subprocess execution sandbox terminated process at 30.0s timeout; UI remained responsive. & \textbf{ISOLATED} \\ \hline
\textbf{Obfuscated Eval} & \texttt{eval(base64.b64decode(...))} & AST analyzer flagged dynamic execution primitive; routed to Tier 3 confirmation prompt. & \textbf{BLOCKED} \\ \hline
\textbf{Unauthorized Mobile PIN} & 100 rapid brute-force PIN guesses on LAN HTTP & HTTP 401 Unauthorized returned; client IP rate-limited after 5 failed attempts. & \textbf{LOCKED} \\ \hline
\textbf{Telegram Spoofing} & Commands sent from un-paired external Telegram chat & Chat ID checked against authenticated constant; update dropped silently without processing. & \textbf{REJECTED} \\ \hline
\textbf{Microphone Sniffing} & Background audio snooping attempt during idle & Audio stream handle confirmed \texttt{NULL} in memory; 0 samples captured during idle state. & \textbf{VERIFIED} \\ \hline
\end{tabularx}
\end{table}

All seven penetration vectors were successfully defended, confirming that PRIVACY68 enforces a robust, multi-layered security posture across all operational modes.
